\documentclass{beamer}
\usepackage{amsmath}
\usepackage{braket}
\usepackage{color}
\usepackage{bbold}
\usepackage{graphicx} % Required for inserting images
\usepackage{ amssymb } % for check marks
\usepackage{svg}
\usepackage{tikz}

\usetikzlibrary{quantikz2}

\usetheme{Antibes}

\title{Steane Type Error Correction - Error Propagation}
\subtitle{Short Long Update}
\author{Leonard Lux}
\date{24.02.2026}

\begin{document}

\tikzset{ noisy/.style={fill=red!20}
}

% Title page
\begin{frame}
\titlepage
\begin{center}
Supervised by:\newline 
\end{center}
\begin{center}
Luis Colmenarez, Maarten Wegewijs \& Markus Müller
\end{center}
\end{frame}



\section{Surface Code}
\begin{frame}{Encoding: \textbf{Surface Code}}
    \begin{columns}
    \begin{column}{0.5\textwidth}
        \begin{itemize}
            \setlength{\itemsep}{8pt}
            \item Stabilizer code
            \item Transversal $\overline{\mathrm{CNOT}}$ 
            \item $[[d^2 + (d-1)^2,1,d]]$
        \end{itemize}
    \end{column}
    \begin{column}{0.5\textwidth}  
        \begin{center}
        \includegraphics[width=0.9\textwidth]{sketch_surface_code.png}
        \end{center}
    \end{column}
    \end{columns}
\end{frame}

\section{Efficient MLD for Surface Code}
\begin{frame}{\textbf{Efficient} MLD for Surface Code}
    \begin{center}
        \includegraphics[width=0.7\textwidth]{eff_ml_decoder_paper.png}
    \end{center}
    \begin{itemize}
        \setlength{\itemsep}{6pt}
        \item Designed for surface codes 
        \item Assumes* independent bit- \& phase-flip noise 
   \end{itemize}
\end{frame}

\section{Steane Type Error Correction}
\begin{frame}{Steane Type Error Correction}
    \begin{columns}
    \begin{column}{0.6\textwidth}
    \begin{block}{The Circuit}
        % Simple Steane Type EC
        \begin{quantikz}
            \lstick{$\ket{\Psi}_L$}&\qwbundle{n}&\targ{}&&\ctrl{2}&&\\%&&\gate{Z_R}&\gate{X_R}&\\
            \lstick{$\ket{0}_L$}&\qwbundle{n}&\ctrl{-1}&\meter{X}&\setwiretype{c}&&\\%\gate{Decoder}\wire[u][1]{c}\\
            \lstick{$\ket{+}_L$}&\qwbundle{n}&&&\targ{}&\meter{Z}&\setwiretype{c}%&\gate{Decoder}\wire[u][2]{c}
        \end{quantikz}
    \end{block} 
    \end{column}
    \begin{column}{0.4\textwidth}
    \begin{itemize}
        \item Fault-tolerant  
        \item Single-shot 
        \item Not symmetric 
    \end{itemize} 
    \end{column}    
    \end{columns}
\end{frame}


% Slide  (Error Models)

\section{Error Models}
\begin{frame}{Circuit level noise model}
    Probability of noise is of $\mathcal{O}(p)$
    \begin{itemize}
        \setlength{\itemsep}{6pt}
        \item Depolarizing single qubit noise ($D1$) 
        \begin{itemize}
            \item After logical encoding of qubit (Assumption!) 
        \end{itemize}
        \item Depolarizing two qubit noise ($D2$) 
        \begin{itemize}
            \item After $\mathrm{CNOT}$ gates
        \end{itemize}
        \item Measurement errors
        \begin{itemize}
            \item Bit-flip before $Z$-Measurements
            \item Phase-flip before $X$-Measurements
        \end{itemize}
    \end{itemize}
\end{frame}


% Slide: Circuit examples Errors
\begin{frame}{Circuit Level Noise}

    \begin{quantikz}
        \lstick{$\ket{\Psi_L}$}&\qwbundle{n}&\targ{}&\ctrl{2}&&\\
        \lstick{$\ket{0_L}$}&\qwbundle{n}&\ctrl{-1}&&\meter{$X$}&\setwiretype{c}\\
        \lstick{$\ket{+_L}$}&\qwbundle{n}&&\targ{}&\meter{$Z$}&\setwiretype{c}\\
    \end{quantikz}

    \begin{quantikz}
        % Data qubit
        \lstick{$\ket{\Psi_L}$}&\qwbundle{n}&\gate[1,style={noisy},]{D}&   % 1: init
        \targ{}&\gate[1,style={noisy}]{D2}\wire[d][1]{q}&      % 2: 1. CNOT 
        \ctrl{2}&\gate[1,style={noisy}]{D2}\wire[d][2]{q}&     % 3: 2. CNOT
        % &&                                                      % 4: 
        &&\\                                                    % 5: 
        % 0 ancilla qubit
        \lstick{$\ket{0_L}$}&\qwbundle{n}&\gate[1,style={noisy},]{D}&      % 1: init
        \ctrl{-1}&\gate[1,style={noisy}]{D2}&                  % 2: 1. CNOT
        &&                                                      % 3: 2. CNOT
        % \gate[]{H}&                                             % 4: 1. change basis aka H
        % & 
        % \gate[1,style={noisy},]{D_1}&                           %    2: Error (No errors on H)
        \gate[1,style={noisy},]{Z}&\meter{$X$}&\setwiretype{c}\\% 5: stabilizer measurement
        % p ancilla qubit 
        \lstick{$\ket{+_L}$}&\qwbundle{n}&\gate[1,style={noisy},]{D}&      % 1: init
        &&                                                      % 2: 1. CNOT
        \targ{-2}&\gate[1,style={noisy}]{D2}&                  % 3: 2. CNOT
        % &&                                                      % 4: 
        \gate[1,style={noisy},]{X}&\meter{$Z$}&\setwiretype{c}\\% 5: stabilizer measurement
    \end{quantikz}
\end{frame}


\subsection{Detector Error Model} 
\begin{frame}{Detector Error Model}
    \begin{quantikz}
        % Data qubit
        \lstick{$\ket{\Psi_L}$}&
        \targ{}&
        \ctrl{2}&
        \meter{trace}
        % &                                                      
        \\                                                    
        % 0 ancilla qubit
        \lstick{$\ket{0_L}$}&
        \ctrl{-1}&
        &
        &
        \gate[1,style={noisy}]{D'}\wire[d][1]{q}&
        \gate[1,style={noisy},]{Z}&\meter{$X$}&\setwiretype{c}\\% 5: stabilizer measurement
        % p ancilla qubit 
        \lstick{$\ket{+_L}$}&
        &
        \targ{-2}&
        &
        \gate[1,style={noisy}]{D'}&
        % &&                                                      % 4: 
        \gate[1,style={noisy},]{X}&\meter{$Z$}&\setwiretype{c}\\% 5: stabilizer measurement
    \end{quantikz}
    \begin{align*}
        &p^c_x \left(\frac{2}{3}p_{sp.\ket{+}}, \frac{2}{3}\frac{4}{5}p_{D2\ket{+}}, p_{m\ket{+}} \right) \\
        &p^c_z \left(p_{m\ket{0}} \right) \\
        &p^c_{D'} \left(p_{sp.\ket{\Psi}},p_{sp.\ket{0}}, \frac{4}{5}p_{D2\ket{0}} \right)  
    \end{align*}
\end{frame}

\begin{frame}{Detector Error Model}
    \begin{block}{Resulting Probabilities (correlated)}
        \begin{equation*}
            p_{(i,j)}^{(MLD)} = \frac{1}{3}p_{D'^c} + (1- \frac{4}{3}p_{D'}^c)\cdot p_{x,i}^c \cdot p_{z,j}^c
        \end{equation*}
        \vspace{0.5cm}
        \begin{equation*}
            p_{x,i} = 
            \begin{cases}
                p_x & \text{if } i=1  \\
                (1-p_x) & \text{if } i=0\\
            \end{cases}
        \end{equation*}
    \end{block}
    \begin{itemize}
        \item ML-Decoder assumes uncorrelated $X$ and $Z$ errors
    \end{itemize}
\end{frame}

\begin{frame}{Detector Error Model (uncorrelated)}
    \begin{itemize}
        \item Look at independent probability of $X-$/$Z-$ Errors
        \begin{itemize}
            \item trace out other ancilla 
            \item $\rightarrow$ composition with combined depolarizing channel
        \end{itemize}
    \end{itemize}
    \begin{align*}
        p^{c'}_x \left(p^{c}_x, \frac{2}{3}p^{c}_D\right)  \\
        p^{c'}_z \left(p^{c}_z, \frac{2}{3}p^{c}_D\right) 
    \end{align*}
    \begin{itemize}
        \item Used in current eff. ML implementation 
    \end{itemize}
\end{frame}

\section{Determine Threshold}
\begin{frame}{Determine Threshold - Finite Size Scaling}
    \begin{itemize}
        \item Goal: determine threshold $p_{th}$
        \item Use: finite size scaling ansatz
        \begin{itemize}
            \item For first order phase transition
            \item Assuming $d \rightarrow \infty$
        \end{itemize}
    \end{itemize}
    \begin{block}{Data Collapse}
    \begin{equation*}
        p_{log,d}(p_{phy}) = f\left( d^{1/\nu} \cdot(p_{phy}-p_{th})\right)
    \end{equation*}
with scaling function $f(x)$ 
    \end{block}
    \begin{itemize}
        \item Result: determine $p_{th}$ (and $\nu$) 
    \end{itemize}
\end{frame}

\section{Results}
\subsection{Code Capacity}
\begin{frame}{Code Capacity Noise Model}
    show plots and results 
        ml and mwpm
    
    compare to literature

    discuss finite size effects
\end{frame}

\subsection{Circuit Level}
\begin{frame}{Circuit Level}
    show results for one round
    discuss effect of ordering 
\end{frame}

\subsection{Multiple Rounds QEC}
\begin{frame}{Multiple Rounds}
    Multiple rounds 
\end{frame}


\begin{frame}
    \begin{center}
        Thanks for listening!
    \end{center}  
\end{frame}

\end{document}